\documentclass{article}
\usepackage[utf8]{inputenc}
\usepackage[T1]{fontenc}
\usepackage{amsmath}
\usepackage{amsfonts}
\usepackage{amssymb}
\usepackage{enumitem}

\usepackage[margin=0.7in]{geometry}

\usepackage{tikz}
\usetikzlibrary{matrix,arrows,decorations.pathmorphing,shapes.geometric,calc,snakes,backgrounds,arrows.meta}
\usepackage{xcolor}

\newcommand{\R}{\mathbb{R}}
\newcommand{\Q}{\mathbb{Q}}
\newcommand{\Z}{\mathbb{Z}}

\begin{document}
\pagestyle{plain}

\section*{Solutions: Selected IMC Problems 2001--2006}

\begin{center}
April 22, 2026
\end{center}

\textbf{Problem 1.} (IMC 2001, Day 2, Problem 3)
Find the maximum number of points on a sphere of radius \( 1 \) in \( \R^n \) such that the distance between any two of these points is strictly greater than \( \sqrt{2} \).

\textbf{Solution.}
For two unit vectors \( X, Y \in S^{n-1} \),
\[
|X - Y| > \sqrt{2} \iff \sum_{k=1}^n (x_k - y_k)^2 > 2 \iff \sum_{k=1}^n x_k y_k < 0.
\]

Let \( a_n \) be the desired maximum. After a rotation we may place one point at \( A_1 = (-1, 0, \ldots, 0) \). For every other point \( Y = (y_1, \overline{Y}) \) the condition \( \langle A_1, Y \rangle < 0 \) forces \( y_1 > 0 \); in particular \( |\overline{Y}|^2 < 1 \). For two such points \( X, Y \) we have \( x_1 y_1 > 0 \), so
\[
\sum_{k=1}^n x_k y_k < 0 \implies \sum_{k=2}^n x_k y_k < 0.
\]
Normalising \( \overline{X}, \overline{Y} \) to unit vectors yields \( a_n - 1 \) points on \( S^{n-2} \) with pairwise distance \( > \sqrt{2} \). Hence \( a_n \leq 1 + a_{n-1} \); since \( a_1 = 2 \) (two antipodal points), \( a_n \leq n + 1 \).

Equality is attained by the \( n+1 \) vertices of a regular simplex inscribed in \( S^{n-1} \): their pairwise inner products equal \( -1/n \), so pairwise distances equal \( \sqrt{2}\sqrt{1 + 1/n} > \sqrt{2} \).

\textit{Answer:} \( a_n = n + 1 \). \hfill\(\square\)
\\

\textbf{Problem 2.} (IMC 2005, Day 2, Problem 2)
Let \( f : \R \to \R \) be a function such that \( (f(x))^n \) is a polynomial for every integer \( n \geq 2 \). Does it follow that \( f \) itself is a polynomial?

\textbf{Solution.}
\textit{Yes.} It suffices that \( f^2 \) and \( f^3 \) are polynomials. Put \( p = f^2 \) and \( q = f^3 \). If \( p \equiv 0 \) then \( f \equiv 0 \), done. Otherwise \( q^2 = f^6 = p^3 \not\equiv 0 \). Define the rational function
\[
r(x) = \frac{q(x)}{p(x)}.
\]
Wherever \( f(x) \neq 0 \), \( r(x) = f(x)^3 / f(x)^2 = f(x) \).

Write \( r = P/Q \) in lowest terms (\( \gcd(P, Q) = 1 \)). Then \( r^2 = P^2/Q^2 \) is also in lowest terms. But
\[
r^2 = \frac{q^2}{p^2} = \frac{p^3}{p^2} = p,
\]
a polynomial. A rational function in lowest terms is a polynomial iff its denominator is constant, so \( Q \) is constant and \( r \) is a polynomial.

Finally, \( f \) and \( r \) agree except at zeros of \( p \). At any zero \( x_0 \) of \( p \), \( f(x_0)^2 = p(x_0) = 0 \) and \( r(x_0)^2 = p(x_0) = 0 \), so \( f(x_0) = r(x_0) = 0 \). Thus \( f = r \) everywhere. \hfill\(\square\)
\\

\textbf{Problem 3.} (IMC 2003, Day 1, Problem 4)
Determine the set of all pairs \( (a, b) \) of positive integers for which the set of positive integers can be decomposed into two sets \( A, B \) such that \( a \cdot A = b \cdot B \).

\textbf{Solution.}
Clearly \( a \neq b \) (else \( A = B \), but \( A \cap B = \emptyset \)). Setting \( d = \gcd(a, b) \), \( a = d a_1 \), \( b = d b_1 \) with \( \gcd(a_1, b_1) = 1 \), the condition \( a \cdot A = b \cdot B \) is equivalent to \( a_1 \cdot A = b_1 \cdot B \). So WLOG \( \gcd(a, b) = 1 \).

If \( 1 \in A \), then \( a \in a \cdot A = b \cdot B \), so \( b \mid a \), hence \( b = 1 \). Symmetrically, if \( 1 \in B \), then \( a = 1 \). So one of \( a, b \) is \( 1 \).

Conversely, for \( (a, b) = (1, n) \) with \( n \geq 2 \), let \( v_n(k) \) denote the largest integer \( e \geq 0 \) with \( n^e \mid k \) (the \( n \)-adic valuation). Define
\[
A = \{k \in \Z_{>0} : v_n(k) \text{ odd}\}, \qquad B = \{k \in \Z_{>0} : v_n(k) \text{ even}\}.
\]
Then \( A \sqcup B = \Z_{>0} \), and multiplying by \( n \) increases \( v_n \) by \( 1 \), so \( n \cdot B = A = 1 \cdot A \).

\textit{Answer:} the pairs \( (a, b) \) with \( a \neq b \) and \( a \mid b \) or \( b \mid a \). \hfill\(\square\)
\\

\textbf{Problem 4.} (IMC 2003, Day 1, Problem 6)
If all roots of \( f(z) = a_n z^n + \cdots + a_0 \) (real coefficients) lie in the left half-plane \( \operatorname{Re} z < 0 \), then \( a_k a_{k+3} < a_{k+1} a_{k+2} \) for \( k = 0, \ldots, n-3 \).

\textbf{Solution.}
Factor \( f \) over \( \R \) as a product of linear and quadratic factors. Each linear factor \( kz + l \) with a root \( -l/k \) in the open left half-plane has \( k, l \) of the same sign; each quadratic factor \( pz^2 + qz + r \) with complex-conjugate roots in the left half-plane has \( p, q, r \) of the same sign. Multiplying \( f \) by \( -1 \) if needed, we may assume \textbf{all coefficients \( a_k \) are positive}.

Extend with \( a_j = 0 \) for \( j < 0 \) or \( j > n \), and prove the stronger claim
\[
a_k a_{k+3} \leq a_{k+1} a_{k+2} \text{ for all } k \in \Z, \quad \text{with strict inequality for } 0 \leq k \leq n-3.
\]
Induct on \( n \). The cases \( n \leq 2 \) are trivial.

For \( n \geq 3 \), choose a real factor \( z^2 + pz + q \) (\( p, q > 0 \)) and write \( f(z) = (z^2 + pz + q) g(z) \) with \( g \) of degree \( n-2 \) with coefficients \( b_0, \ldots, b_{n-2} > 0 \). Then \( a_k = q b_k + p b_{k-1} + b_{k-2} \), and a direct expansion gives
\begin{align*}
a_{k+1} a_{k+2} - a_k a_{k+3}
&= \bigl(b_{k-1} b_k - b_{k-2} b_{k+1}\bigr)
+ p \bigl(b_k^2 - b_{k-2} b_{k+2}\bigr)
+ q \bigl(b_{k-1} b_{k+2} - b_{k-2} b_{k+3}\bigr) \\
&\quad + p^2 \bigl(b_k b_{k+1} - b_{k-1} b_{k+2}\bigr)
+ q^2 \bigl(b_{k+1} b_{k+2} - b_k b_{k+3}\bigr)
+ pq \bigl(b_{k+1}^2 - b_{k-1} b_{k+3}\bigr).
\end{align*}
Each of the six bracketed expressions is non-negative (using the inductive hypothesis for \( g \); for the square terms like \( b_k^2 - b_{k-2} b_{k+2} \), use the identity
\( b_{k-1}(b_k^2 - b_{k-2} b_{k+2}) = b_{k-2}(b_k b_{k+1} - b_{k-1} b_{k+2}) + b_k(b_{k-1} b_k - b_{k-2} b_{k+1}) \geq 0 \)
and handle \( b_{k-1} = 0 \) separately). For \( 0 \leq k \leq n-3 \) the term \( p^2 (b_k b_{k+1} - b_{k-1} b_{k+2}) \) is strictly positive, so the inequality is strict. \hfill\(\square\)
\\

\textbf{Problem 5.} (IMC 2004, Day 2, Problem 6)
With \( A_n, B_n \) as defined and \( S(M) \) the sum of entries, \( S(A_n^{k-1}) = S(A_k^{n-1}) \).

\textbf{Solution.}
Let \( F_{n,k} \) be the number of \( n \times k \) \( \{0,1\} \)-matrices with no \( 2 \times 2 \) all-ones sub-square. Transposition gives \( F_{n,k} = F_{k,n} \). We show \( F_{n,k} = S(A_n^{k-1}) \), which finishes the proof.

Encode each column (a \( 0/1 \)-vector of length \( n \)) as an integer \( 0 \leq i < 2^n \) with binary digits \( i_n \ldots i_1 \). Two consecutive columns \( i, j \) are \emph{compatible} (avoid a forbidden \( 2 \times 2 \) block) iff for every adjacent pair of rows \( \ell, \ell - 1 \) we do not have \( i_\ell = i_{\ell-1} = j_\ell = j_{\ell-1} = 1 \). Let \( C^{(n)} \) be the \( 2^n \times 2^n \) compatibility matrix. Then \( F_{n,k} = S\bigl((C^{(n)})^{k-1}\bigr) \) (enumerate all sequences of \( k \) columns).

We show \( C^{(n)} = A_n \) by induction on \( n \). Block indices by the top bits \( i_n, j_n \):
\begin{itemize}
    \item \( i_n = 0 \) or \( j_n = 0 \): compatibility of columns reduces to compatibility of the lower \( n-1 \) bits, giving blocks equal to \( C^{(n-1)} = A_{n-1} \). This is the top-left, top-right, and bottom-left blocks of \( A_n \).
    \item \( i_n = j_n = 1 \): we additionally require \( i_{n-1} j_{n-1} = 0 \), and compatibility of the lower \( n-2 \) bits. This is exactly \( B_{n-1} \) (the bottom-right block of \( A_n \)): starts like \( A_{n-1} \) but the bottom-right \( 2^{n-2} \times 2^{n-2} \) sub-block is zeroed out, since that would correspond to \( i_{n-1} = j_{n-1} = 1 \) too.
\end{itemize}
Thus \( C^{(n)} = A_n \), and \( S(A_n^{k-1}) = F_{n,k} = F_{k,n} = S(A_k^{n-1}) \). \hfill\(\square\)
\\

\textbf{Problem 6.} (IMC 2005, Day 1, Problem 4)
Find all polynomials \( P(x) = a_n x^n + \cdots + a_0 \) with \( a_n \neq 0 \) such that \( (a_0, \ldots, a_n) \) is a permutation of \( (0, \ldots, n) \) and all roots are rational.

\textbf{Solution.}
All coefficients are non-negative, so \( P(x) > 0 \) for \( x > 0 \); all real roots are \( \leq 0 \). Rational roots force a factorisation over \( \Z \):
\[
P(x) = \prod_{k=1}^n (b_k x + c_k), \qquad b_k > 0, \; c_k \geq 0.
\]
At most one \( c_k \) is \( 0 \) (else both \( a_0 \) and \( a_1 \) vanish), so at least \( n-1 \) of the factors have \( b_k + c_k \geq 2 \).

Evaluating at \( x = 1 \):
\[
P(1) = \sum_{k=0}^n a_k = 0 + 1 + \cdots + n = \frac{n(n+1)}{2} = \prod_{k=1}^n (b_k + c_k) \geq 2^{n-1}.
\]
Hence \( 2^{n-1} \leq n(n+1)/2 \), forcing \( n \leq 4 \). Moreover \( n(n+1)/2 \) must factor as a product of \( n-1 \) integers \( \geq 2 \):
\begin{itemize}
    \item \( n = 1 \): \( P(x) = x \).
    \item \( n = 2 \): \( P(1) = 3 = 1 \cdot 3 \); solutions \( x(x+2) = x^2 + 2x \) and \( x(2x+1) = 2x^2 + x \).
    \item \( n = 3 \): \( P(1) = 6 = 1 \cdot 2 \cdot 3 \); solutions \( x(x+1)(x+2) = x^3 + 3x^2 + 2x \) and \( x(x+1)(2x+1) = 2x^3 + 3x^2 + x \).
    \item \( n = 4 \): \( P(1) = 10 \) cannot be written as a product of 3 integers each \( \geq 2 \). No solutions.
\end{itemize}

\textit{Answer:} \( P(x) \in \{x,\; x^2+2x,\; 2x^2+x,\; x^3+3x^2+2x,\; 2x^3+3x^2+x\} \). \hfill\(\square\)
\\

\textbf{Problem 7.} (IMC 2005, Day 2, Problem 1)
\( |M| \leq 2\sqrt{2} \), where \( M = \{x \in \R : |x^2+bx+c| < 1\} \).

\textbf{Solution.}
Complete the square: \( f(x) = (x + b/2)^2 + d \) where \( d = c - b^2/4 \). With \( u = x + b/2 \):
\begin{itemize}
    \item \( d \geq 1 \): \( f \geq 1 \), so \( M = \emptyset \), \( |M| = 0 \).
    \item \( -1 < d < 1 \): \( |f| < 1 \iff u^2 < 1 - d \), so \( |M| = 2\sqrt{1-d} < 2\sqrt{2} \).
    \item \( d \leq -1 \): \( |f| < 1 \iff \sqrt{|d|-1} < |u| < \sqrt{|d|+1} \), giving
    \[
    |M| = 2\bigl(\sqrt{|d|+1} - \sqrt{|d|-1}\bigr) = \frac{4}{\sqrt{|d|+1} + \sqrt{|d|-1}}.
    \]
    This is a decreasing function of \( |d| \geq 1 \), maximised at \( |d| = 1 \): \( |M|_{\max} = \frac{4}{\sqrt{2} + 0} = 2\sqrt{2} \).
\end{itemize}
In all cases \( |M| \leq 2\sqrt{2} \). \hfill\(\square\)
\\

\textbf{Problem 8.} (IMC 2006, Day 1, Problem 1)
Prove or disprove: (a) continuous + surjective \( \Rightarrow \) monotonic; (b) monotonic + surjective \( \Rightarrow \) continuous; (c) monotonic + continuous \( \Rightarrow \) surjective.

\textbf{Solution.}

\textit{(a) False.} Take \( f(x) = x^3 - x \): continuous and \( \lim_{x \to \pm\infty} f(x) = \pm\infty \) gives surjectivity, but \( f(0) = 0 \), \( f(1/2) = -3/8 \), \( f(1) = 0 \), so \( f \) is not monotonic.

\textit{(b) True.} Assume (WLOG) \( f \) non-decreasing. For \( a \in \R \), the one-sided limits \( L^- = \lim_{x \to a^-} f(x) \) and \( L^+ = \lim_{x \to a^+} f(x) \) exist in \( [-\infty, +\infty] \) with \( L^- \leq f(a) \leq L^+ \). If \( L^- < L^+ \), then \( \operatorname{range}(f) \subseteq (-\infty, L^-] \cup \{f(a)\} \cup [L^+, +\infty) \), missing an open interval — contradicting surjectivity. So \( L^- = L^+ = f(a) \), and \( f \) is continuous at \( a \).

\textit{(c) False.} \( g(x) = \arctan x \) is continuous and strictly increasing, but \( \operatorname{range}(g) = (-\pi/2, \pi/2) \neq \R \). \hfill\(\square\)
\\

\textbf{Problem 9.} (IMC 2006, Day 1, Problem 5)
With \( a, b, c, d, e > 0 \), \( a^2+b^2+c^2 = d^2+e^2 \) and \( a^4+b^4+c^4 = d^4+e^4 \), compare \( a^3+b^3+c^3 \) with \( d^3+e^3 \).

\textbf{Solution.}
\textit{Answer:} \( a^3 + b^3 + c^3 < d^3 + e^3 \).

WLOG \( a \geq b \geq c > 0 \) and \( d \geq e > 0 \). Squaring \( a^2 + b^2 + c^2 = d^2 + e^2 \) and using \( a^4 + b^4 + c^4 = d^4 + e^4 \):
\[
a^2 b^2 + a^2 c^2 + b^2 c^2 = d^2 e^2.
\]
With \( x = a^2, y = b^2, z = c^2, p = d^2, q = e^2 \), the polynomials \( (t-x)(t-y)(t-z) \) and \( t(t-p)(t-q) \) agree in degrees \( 3, 2, 1 \), so
\[
(t-x)(t-y)(t-z) = t(t-p)(t-q) - xyz. \qquad (\star)
\]

\textit{Ordering.} Evaluate \( (\star) \) at \( t = p \): \( (p-x)(p-y)(p-z) = -xyz < 0 \). If all three factors were negative, then \( x + y + z > 3p \geq p + q \), contradicting \( x + y + z = p + q \). So exactly one factor is negative; since \( x \geq y \geq z \), this forces \( x > p \). Then \( y + z = (p+q) - x < q \leq p \), so \( y, z < q \leq p \). Hence
\[
a > d \geq e > b \geq c.
\]

\textit{Comparison.} Let \( F(r) = a^r + b^r + c^r - d^r - e^r \) for \( r \in \R \). Then \( F(0) = 1 \), \( F(2) = F(4) = 0 \); we want the sign of \( F(3) \).

We claim \( F \) has no zero in \( (0, \infty) \) other than \( r = 2 \) and \( r = 4 \). Writing \( F(r) = \sum_{i=1}^5 \varepsilon_i \alpha_i^r \) with positive bases \( \alpha_1 = a > \alpha_2 = d \geq \alpha_3 = e > \alpha_4 = b \geq \alpha_5 = c \) and signs \( (+, -, -, +, +) \), the sequence of signs has exactly \textbf{2 sign changes} (at positions \( 1 \to 2 \) and \( 3 \to 4 \)). By a generalized Descartes' rule of signs for exponential polynomials (``Laguerre's theorem''), \( F(r) \) has at most \( 2 \) positive real zeros, counted with multiplicity. So \( r = 2, 4 \) are the only zeros.

Since \( F(0) = 1 > 0 \), \( F \) is positive on \( (-\infty, 2) \); by continuity and the absence of other zeros, \( F \) is negative on \( (2, 4) \) and positive on \( (4, \infty) \). In particular,
\[
F(3) = (a^3 + b^3 + c^3) - (d^3 + e^3) < 0. \qquad\square
\]
\\

\textbf{Problem 10.} (IMC 2006, Day 2, Problem 4)
With rational \( |v_i - v_j| \) for \( 0 \leq i, j \leq n+1 \) and \( v_0 = 0 \), vectors \( v_1, \ldots, v_{n+1} \in \R^n \) are \( \Q \)-linearly dependent.

\textbf{Solution.}
If \( v_1, \ldots, v_n \) are already \( \R \)-linearly dependent, they (hence together with \( v_{n+1} \)) are \( \Q \)-linearly dependent by clearing a non-zero real relation to a rational one via a sub-matrix. So assume \( v_1, \ldots, v_n \) are \( \R \)-linearly independent, and write
\[
v_{n+1} = \sum_{j=1}^n \lambda_j v_j, \qquad \lambda_j \in \R.
\]
We show \( \lambda_j \in \Q \).

By polarization, using \( v_0 = 0 \):
\[
\langle v_i, v_j \rangle = \tfrac{1}{2}\bigl( |v_i|^2 + |v_j|^2 - |v_i - v_j|^2 \bigr) \in \Q.
\]
The Gram matrix \( A = [\langle v_i, v_j \rangle]_{i,j=1}^n \in \Q^{n \times n} \) is non-singular (since \( v_1, \ldots, v_n \) are linearly independent) and the vector \( w = (\langle v_i, v_{n+1} \rangle)_{i=1}^n \in \Q^n \).

Taking \( \langle \cdot, v_i \rangle \) of both sides of \( v_{n+1} = \sum \lambda_j v_j \):
\[
A \lambda = w \implies \lambda = A^{-1} w \in \Q^n.
\]
So \( \lambda_j \in \Q \), giving a non-trivial \( \Q \)-linear relation among \( v_1, \ldots, v_{n+1} \). \hfill\(\square\)

\end{document}
